\noindent\textbf{Einheit 2 · Prozentsatz berechnen}

\erg{8}{a) die ganze Klasse, 25 Kinder \quad b) alle 60 Zuschauer \quad c) der alte Preis, $400\,€$ \quad d) alle Äpfel zusammen, 20 \quad e) alle 30 Lose}
\erg{9}{a) $30\,\%$ \quad b) $70\,\%$ \quad c) $40\,\%$ \quad d) $45\,\%$}
\erg{10}{a) $19\,\%$ \quad b) $7\,\%$ \quad c) $64\,\%$ \quad d) $90\,\%$ \quad e) $84\,\%$ \quad f) $35\,\%$ \quad g) $5 : 12 = 0{,}4166\ldots \approx 41{,}7\,\%$}
\erg{11}{a) z.\,B. Teil 8, Ganzes 16 \quad b) z.\,B. Teil 5, Ganzes 20 \quad c) z.\,B. 2 und 5 sowie 20 und 50 (jedes Paar mit Teil $:$ Ganzes $= 0{,}4$)}
\erg{12}{a) $26 : 40 = 0{,}65 = 65\,\%$ \quad b) Ganzes 50; $18 : 50 = 36\,\%$ \quad c) Ersparnis $12\,€$; $12 : 80 = 15\,\%$ \quad d) Ganzes 24; $7 : 24 = 0{,}2916\ldots \approx 29{,}2\,\%$}
\erg{13}{a) Nina hat das Ganze durch den Teil geteilt; richtig ist $24 : 40 = 0{,}6 = 60\,\%$ \quad b) $36 : 45 = 0{,}8 = 80\,\%$}
\erg{14}{a) Der Prozentsatz sagt, welcher Anteil vom Ganzen der Teil ist; das Ganze ist $100\,\%$ und steht deshalb im Nenner \quad b) ja – ein Teil kann nie größer sein als das Ganze, aus dem er genommen wird}
\erg{15}{a) $12 : 20 = 60\,\%$ \quad b) $21 : 35 = 60\,\%$ \quad c) beide gleich gut; Ben trifft öfter, wirft aber auch öfter}
